\documentclass[a4paper]{arrowhead}

\usepackage[yyyymmdd]{datetime}
\usepackage{etoolbox}
\usepackage[utf8]{inputenc}
\usepackage[T1]{fontenc}
\usepackage{multirow}
\usepackage{longtable}
\usepackage{tikz}
\usetikzlibrary{positioning,shapes.multipart,arrows.meta}

\renewcommand{\dateseparator}{-}
\setlength{\emergencystretch}{3em}
\lstset{language=HTTP,breakatwhitespace=false,columns=fullflexible}

%% Special references
\newcommand{\mref}[2][]{{\textcolor{ArrowheadPurple}{\hyperref[sec:model:#2]{\ifstrempty{#1}{#2}{#1}}}}}
\newcommand{\oref}[1]{{\textcolor{ArrowheadBlue}{\hyperref[sec:operations:#1]{#1}}}}
\newcommand{\pdef}[1]{{\textcolor{ArrowheadGrey}{#1\label{sec:model:primitives:#1}\label{sec:model:primitives:#1s}\label{sec:model:primitives:#1es}}}}
\newcommand{\pref}[1]{{\textcolor{ArrowheadGrey}{\hyperref[sec:model:primitives:#1]{#1}}}}

% as template_sd.tex; the input and output arguments are passed already formatted (\mref or \pref)
\newrobustcmd\osubsection[3]{
  \addtocounter{subsection}{1}
  \addcontentsline{toc}{subsection}{\protect\numberline{\thesubsection}operation \textcolor{ArrowheadBlue}{#1}}
  \subsection*{
    \thesubsection\quad
    operation
    \textcolor{ArrowheadBlue}{#1}
    (\notblank{#2}{#2}{})
    \notblank{#3}{: #3}{}
  }
  \label{sec:operations:#1}
}
\newrobustcmd\msubsection[2]{
  \addtocounter{subsection}{1}
  \addcontentsline{toc}{subsection}{\protect\numberline{\thesubsection}#1 \textcolor{ArrowheadPurple}{#2}}
  \subsection*{\thesubsection\quad#1 \textcolor{ArrowheadPurple}{#2}}
  \label{sec:model:#2} \label{sec:model:#2s} \label{sec:model:#2es}
}
%%

\begin{document}
\sloppy

%% Arrowhead Document Properties
\ArrowheadTitle{certificateValidation}
\ArrowheadServiceID{certificateValidation}
\ArrowheadType{Service Description}
\ArrowheadTypeShort{SD}
\ArrowheadVersion{5.2.2}
\ArrowheadDate{2026-10-30}
\ArrowheadAuthor{Jerker Delsing}
\ArrowheadStatus{RELEASE CANDIDATE}
\ArrowheadContact{jerker.delsing@ltu.se}
\ArrowheadFooter{\href{www.arrowhead.eu}{www.arrowhead.eu}}
\ArrowheadSetup
%%

%% Front Page
\begin{center}
  \vspace*{1cm}
  \huge{\arrowtitle}

  \vspace*{0.2cm}
  \LARGE{\arrowtype}
  \vspace*{1cm}

  \Large{Service ID: \textit{"\arrowid"}}

  \vspace*{\fill}

  % Front Page Image

  \vspace*{1cm}
  \vspace*{\fill}

  % Front Page Abstract
  \begin{abstract}
    This document describes the \texttt{certificate\allowbreak{}Validation} service, version 1.0.0, produced by the CertificateAuthority support system of Arrowhead 5.2.2.
  \end{abstract}

  \vspace*{1cm}
\end{center}

\newpage
%%

%% Table of Contents
\tableofcontents
\newpage
%%

\section{Overview}
\label{sec:overview}

This document describes the \texttt{certificate\allowbreak{}Validation} service, which provides checking certificates and trusted keys against the CA's register, and publishing the certificate revocation list.
It is produced by the CertificateAuthority system \cite{cau-sysd} in service version 1.0.0.

The rest of this document is organized as follows.
In Section \ref{sec:operations}, we describe the abstract message operations provided by the service.
In Section \ref{sec:model}, we end the document by presenting the data types used by the mentioned operations.
The concrete interface is specified in the HTTP/TLS/JSON IDD \cite{cau-certificateValidation-http-idd}.

\subsection{Significant Prior Art}

Successor of the Arrowhead 4.6 Certificate Authority validation endpoints; the CRL is new.

\subsection{How This Service Is Meant to Be Used}

Systems check a presented certificate with \texttt{check-\allowbreak{}certificate} (status \texttt{GOOD}, \texttt{REVOKED}, \texttt{EXPIRED} or \texttt{UNKNOWN}) and a public key with \texttt{check-\allowbreak{}trusted-\allowbreak{}key}. Brokers, the relay and the Java systems use the CRL, which is also written to \texttt{CRL\_\allowbreak{}DIR/\allowbreak{}ca.\allowbreak{}crl} at start, after every revocation and periodically.

\subsection{Important Delimitations and Dependencies}
\label{sec:delimitations}

Only certificates of this cloud CA can be \texttt{GOOD}.

\newpage

\section{Service Interface}
\label{sec:operations}

This section describes the interfaces to the \texttt{certificate\allowbreak{}Validation} service.
In particular, each subsection names an abstract operation, an input type and an output type, in that order.
The input type is named inside parentheses, while the output type is preceded by a colon.
Input and output types are only denoted when accepted or returned, respectively, by the interface in question.

All abstract data types named in this section are defined in Section \ref{sec:model}.

\begin{figure}[ht!]
  \centering
  \begin{tikzpicture}[font=\small]
    \node[draw=ArrowheadBlue, very thick, rectangle split, rectangle split parts=2, align=center,
          rectangle split part fill={ArrowheadBlue!12,white}, text width=9cm, inner sep=6pt] (b)
      {\guillemotleft{}Service\guillemotright{}\\\textbf{certificateValidation}
        \nodepart{second} \texttt{check-\allowbreak{}certificate}\\ \texttt{check-\allowbreak{}trusted-\allowbreak{}key}\\ \texttt{get-\allowbreak{}crl}};
  \end{tikzpicture}
  \caption{SysML-style block of the \texttt{certificate\allowbreak{}Validation} service and its operations.}
  \label{fig:SD}
\end{figure}

The following interface operations are available.

\osubsection{check-certificate}{\mref{CertificateCheckRequest}}{\mref{CertificateCheckResult}}

Returns the status of a certificate.

\osubsection{check-trusted-key}{\mref{TrustedKeyCheckRequest}}{\mref{TrustedKeyCheckResult}}

Returns whether a public key is trusted now.

\osubsection{get-crl}{}{}

Returns the current CRL (every revoked, unexpired certificate; next update after \texttt{CRL\_\allowbreak{}VALIDITY\_\allowbreak{}HOURS}, 168).

\newpage

\section{Information Model}
\label{sec:model}

Here, all data objects that can be part of the \texttt{certificate\allowbreak{}Validation} service are listed in alphabetic order.
Note that each subsection, which describes one type of object, begins with the \textit{struct} keyword, which is used to denote a collection of named fields, each with its own data type.
As a complement to the explicitly defined types in this section, there is also a list of implicit primitive types in Section \ref{sec:model:primitives}, which are used to represent things like hashes and identifiers.

\msubsection{struct}{CertificateCheckRequest}




\noindent\begin{tabularx}{\textwidth}{| >{\raggedright\arraybackslash}p{3.9cm} | >{\raggedright\arraybackslash\small}p{5.3cm} | X |} \hline
\rowcolor{gray!33} Field & Type & Description \\ \hline
certificate & \pref{String} & Certificate to check, PEM or Base64 DER (alias \texttt{encodedCertificate}). \\ \hline
\end{tabularx}



\msubsection{struct}{CertificateCheckResult}




\noindent\begin{tabularx}{\textwidth}{| >{\raggedright\arraybackslash}p{3.9cm} | >{\raggedright\arraybackslash\small}p{5.3cm} | X |} \hline
\rowcolor{gray!33} Field & Type & Description \\ \hline
status & \pref{String} & \texttt{GOOD}, \texttt{REVOKED}, \texttt{EXPIRED} or \texttt{UNKNOWN} (not issued by this CA). \\ \hline
valid & \pref{Boolean} & True only for \texttt{GOOD}. \\ \hline
serial\allowbreak{}Number & \pref{String} & Serial number, decimal. \\ \hline
known & \pref{Boolean} & Whether the serial is in the register (only for certificates of this CA). \\ \hline
\end{tabularx}



\msubsection{struct}{TrustedKeyCheckRequest}




\noindent\begin{tabularx}{\textwidth}{| >{\raggedright\arraybackslash}p{3.9cm} | >{\raggedright\arraybackslash\small}p{5.3cm} | X |} \hline
\rowcolor{gray!33} Field & Type & Description \\ \hline
public\allowbreak{}Key & \pref{String} & PEM Subject\allowbreak{}Public\allowbreak{}Key\allowbreak{}Info (alias \texttt{publicKeyPem}). \\ \hline
\end{tabularx}



\msubsection{struct}{TrustedKeyCheckResult}




\noindent\begin{tabularx}{\textwidth}{| >{\raggedright\arraybackslash}p{3.9cm} | >{\raggedright\arraybackslash\small}p{5.3cm} | X |} \hline
\rowcolor{gray!33} Field & Type & Description \\ \hline
trusted & \pref{Boolean} & Whether the key is a trusted key within its validity. \\ \hline
key & \mref{TrustedKeyRecord} & The matching record, or null. \\ \hline
\end{tabularx}



\msubsection{struct}{TrustedKeyRecord}




\noindent\begin{tabularx}{\textwidth}{| >{\raggedright\arraybackslash}p{3.9cm} | >{\raggedright\arraybackslash\small}p{5.3cm} | X |} \hline
\rowcolor{gray!33} Field & Type & Description \\ \hline
id & \pref{Number} & Record identifier. \\ \hline
public\_key\_pem & \pref{String} & The key, PEM. \\ \hline
description & \pref{String} & Description. \\ \hline
valid\_after & \pref{String} & Start of validity. \\ \hline
valid\_before & \pref{String} & End of validity. \\ \hline
created\_at & \pref{String} & Creation time. \\ \hline
\end{tabularx}



\subsection{Primitives}
\label{sec:model:primitives}

Types and structures mentioned throughout this document that are assumed to be available to implementations of this service.
The concrete interpretations of each of these types and structures must be provided by any IDD document claiming to implement this service.

\noindent\begin{tabularx}{\textwidth}{| p{4.6cm} | X |} \hline
\rowcolor{gray!33} Type & Description \\ \hline
\pdef{Boolean} & One out of \texttt{true} or \texttt{false}. \\ \hline
\pdef{Number} & An integer number. \\ \hline
\pdef{String} & A chain of characters. \\ \hline
\end{tabularx}

\newpage

\bibliographystyle{IEEEtran}
\bibliography{bibliography}

\newpage

\section{Revision History}
\subsection{Amendments}

\noindent\begin{tabularx}{\textwidth}{| p{1cm} | p{3cm} | p{2cm} | X | p{4cm} |} \hline
\rowcolor{gray!33} No. & Date & Version & Subject of Amendments & Author \\ \hline

1 & 2026-10-22 & \arrowversion & Document created from the Arrowhead LaTeX templates for the 5.2.2 Java Spring implementation, including the local extensions of this workspace. & Jerker Delsing \\ \hline
2 & 2026-10-30 & \arrowversion & Checked against the primary Java source tree (Eclipse-Arrowhead-java): services, operations, paths and data models unchanged. & Claude Opus 5.5 (draft) \\ \hline
\end{tabularx}

\subsection{Quality Assurance}

\noindent\begin{tabularx}{\textwidth}{| p{1cm} | p{3cm} | p{2cm} | X |} \hline
\rowcolor{gray!33} No. & Date & Version & Approved by \\ \hline

1 & & \arrowversion  &  \\ \hline

\end{tabularx}

\end{document}
